%% -!TEX root = AUTthesis.tex
% در این فایل، عنوان پایان‌نامه، مشخصات خود، متن تقدیمی‌، ستایش، سپاس‌گزاری و چکیده پایان‌نامه را به فارسی، وارد کنید.
% توجه داشته باشید که جدول حاوی مشخصات پروژه/پایان‌نامه/رساله و همچنین، مشخصات داخل آن، به طور خودکار، درج می‌شود.
%%%%%%%%%%%%%%%%%%%%%%%%%%%%%%%%%%%%
% دانشکده، آموزشکده و یا پژوهشکده  خود را وارد کنید
\faculty{دانشکده مهندسی کامپیوتر و فناوری اطلاعات}
% گرایش و گروه آموزشی خود را وارد کنید
\department{معماری سیستم‌های کامپیوتری}
% عنوان پایان‌نامه را وارد کنید
\fatitle{طراحی و پیاده‌سازی سامانه ردیابی مبتنی بر 
	\\[.75 cm]
	اینترنت اشیا}

 
% نام استاد(ان) راهنما را وارد کنید
\firstsupervisor{دکتر بهادر بخشی}
%\secondsupervisor{استاد راهنمای دوم}
% نام استاد(دان) مشاور را وارد کنید. چنانچه استاد مشاور ندارید، دستور پایین را غیرفعال کنید.
\firstadvisor{دکتر مهدی راستی}
%\secondadvisor{استاد مشاور دوم}
% نام نویسنده را وارد کنید
\name{ ساره سلطانی نژاد}
% نام خانوادگی نویسنده را وارد کنید
\surname{}
%%%%%%%%%%%%%%%%%%%%%%%%%%%%%%%%%%
\thesisdate{بهار
	 98}

% چکیده پایان‌نامه را وارد کنید
\fa-abstract{
در علم فناوری اطلاعات، مفهوم اينترنت اشيا به اشيايي با هويت خاص اطلاق مي‌شود كه دارای شناسه منحصر به فرد بوده و توانايي انتقال داده روی شبکه، بدون نياز به تعامل و دخالت انسان را دارند. در واقع هدف اصلي آن هوشمند سازی اشيا و فراهم آوردن بستری است كه از طريق آن،
اشيا قادر به ارسال و دريافت اطلاعات با يکديگر مي‌باشند.
در سال‌های اخير فناوری اينترنت اشيا رشد چشمگيری داشته و در زمينه‌های مختلف توانسته نيازهای متعدد و پيچيده‌ای را برطرف كند. یکی از این زمینه‌ها ردیابی اشخاص و وسایل نقلیه است. 
\\
سيستم موقعيت‌يابي و رديابي امکان ارائه راه‌حل‌هايي مطمئن برای تامين امنيت افراد و وسايل نقليه را فراهم آورده است و همچنين تاثير بسزايي در بهينه شدن كيفيت نظارت و مديريت
ناوگان‌های حمل و نقل، حركت خودروها، افراد (کودکان و سالمندان) و يا هر شي متحرك ديگر دارد. در واقع سامانه رديابي تکنولوژی است كه امکان تعيين موقعيت دقيق و رديابي افراد، وسايل نقليه و يا هر جسم متحرك ديگر را با استفاده از متدهای مختلفي مانند
سامانه موقعيتإياب جهاني فراهم آورده است.
همان‌طور كه گفتيم اينترنت اشيا يک بستر ارتباطي جديد در جهت برقراری ارتباط بين اشيا هوشمند مي‌باشد. معرفي اين بستر موجب شده است تا امکانات جديدی برای حل مسائلي همچون تعيين مکان، رديابي اشخاص و وسايل نقليه فراهم گردد. پس يکي از كاربردهايي كه مي‌توان برای اينترنت اشيا متصور شد، پياده‌سازی سامانه‌ای است كه بتوان توسط آن موقعيت دقيق و مسير حركت هر جسم متحرك را در هر زمان تعيين كرد. در اين پروژه قصد داريم چنين سامانه‌ای را پياده‌سازی كنيم.
}


% کلمات کلیدی پایان‌نامه را وارد کنید
\keywords{موقعیت‌یاب جهانی، اینترنت اشیا، ردیابی}



\AUTtitle
%%%%%%%%%%%%%%%%%%%%%%%%%%%%%%%%%%
\vspace*{7cm}
\thispagestyle{empty}
\begin{center}
\includegraphics[height=5cm,width=12cm]{besm}
\end{center}